\answer{ex:11:1}{(a)~$T=k^2/\bigl(\Phi(\Delta I/I)^2\bigr)=9\,\text{s}$. (b)~$45$ capteurs, tache $\approx7\times7$, résolution $\approx7$ fois moins bonne.}
\answer{ex:11:2}{(a)~$2$--$3.3$ bits par impulsion, $22$--$34\,\%$ de la limite ($9.1$--$9.8$ bits). (b)~$\log_2\binom{400}{20}\approx111$ bits; en moyenne $1$ type commun.}
\answer{ex:11:3}{(a)~$\sigma/9\le I\le9\sigma$, $1.9$ ordre de grandeur. (b)~$6$ et $7$.}
\answer{ex:11:4}{(a)~$f<343/0.44\approx780$~Hz. (b)~$625$ impulsions, $125$ neurones.}
\answer{ex:11:5}{$21$ bits par événement, $4.8\cdot10^5$ événements/s, $7$ événements par pixel du bord: $136$ pixels/s. Caméra ordinaire: $1.25$ image/s.}
\answer{ex:11:6}{En logarithmes, $\approx1.0\,\text{s}$ dans les deux sens (théorie $0.96\tau$); linéairement, $0.2\,\text{s}$ vers le haut et $3.0\,\text{s}$ vers le bas (théorie $0.18\tau$ et $3.0\tau$).}
\answer{ex:11:7}{(a)~$\ln I$ suit une loi logistique de médiane $\ln\sigma$ et d'échelle $1$: écart $\pi/\sqrt3$ en logarithmes naturels, $0.79$ ordre de grandeur. (b)~$r=r_{\max}\Phi_I(I)^2$.}
